\begin{abstract}
Scientific datasets often contain subtle anomalies caused by sensor faults, experimental drift, measurement noise, rare events, or abnormal signal behavior. In many scientific workflows, anomaly detection must be not only accurate but also interpretable, because researchers need to understand whether unusual observations reflect meaningful phenomena, instrument error, or changes in experimental conditions. This paper presents a SciPy-based anomaly detection workflow using statistical and signal-processing tools. Rather than evaluating the workflow only on a single benchmark dataset, we construct a controlled multi-domain experiment using six scientific signal domains: EEG-like signals, environmental sensor signals, spectroscopy-like signals, vibration signals, astronomical photometry-like signals, and fluid-sensor signals. Five controlled anomaly mechanisms are injected into these signals: spike, drift, dropout, oscillation disturbance, and distribution shift. SciPy statistical, signal-processing, residual, spectral, local-window, and complexity features are extracted from every signal. The experiment evaluates binary anomaly detection, anomaly-mechanism classification, relative feature-group signatures, and rule-based diagnosis. In the controlled experiment, SciPy-derived features achieved perfect binary anomaly detection, but anomaly-mechanism classification was substantially harder, with the best SciPy-feature model achieving macro-F1 = 0.4578 across five anomaly types. A hand-written rule-based diagnosis system using relative feature-group signatures achieved only 0.1993 accuracy, approximately chance. This shows that detection performance does not imply diagnostic interpretability. To test whether the workflow transfers beyond controlled injection, we also evaluate it on the UCR Time Series Anomaly Archive \citep{wu2021ucr}. Because the controlled benchmark provides known anomaly labels while the UCR archive provides only normal training regions, the two experiments evaluate complementary supervised and unsupervised settings. Across 56 evaluated real anomaly files, the best SciPy method achieved mean point-adjusted F1 = 0.2884, ROC-AUC = 0.7758, PR-AUC = 0.1801, and event recall = 0.4107. These results show that SciPy features provide a useful and transparent anomaly-detection baseline, but also reveal the limits of global feature summaries: real anomaly detection remains difficult, and mechanism diagnosis requires more than simple feature-group explanations.
\end{abstract}


\section{Introduction}

Scientific data is rarely perfectly clean. Measurements collected from sensors, laboratory instruments, simulations, biological systems, telescopes, and physical experiments may contain unexpected behavior caused by noise, sensor faults, experimental drift, environmental interference, or rare events. Detecting these anomalies is important because abnormal observations can distort downstream analysis, reduce model reliability, and lead to incorrect scientific conclusions.

However, anomaly detection in scientific data is not only a prediction problem. It is also a diagnostic problem. A researcher often needs to know why a signal is unusual. For example, an abnormal measurement may be caused by a transient spike, a gradual drift, a dropout or flatline, a spectral disturbance, or a distribution shift. A model that only reports ``anomaly'' may be useful for flagging data, but it may not provide enough information to support scientific interpretation or corrective action.

This paper studies anomaly detection using SciPy's statistical and signal-processing tools. SciPy is a foundational scientific computing library for Python that provides algorithms for statistics, signal processing, optimization, integration, interpolation, linear algebra, and other scientific computing tasks \citep{virtanen2020scipy}. The \texttt{scipy.stats} module provides probability distributions, summary statistics, and statistical tests. The \texttt{scipy.signal} module provides tools for filtering, peak detection, spectral analysis, and time-series signal processing. These tools make SciPy a practical foundation for anomaly detection workflows that need to be lightweight, reproducible, and interpretable.

The original motivation for this work is that scientific datasets often contain subtle anomalies caused by sensor faults, experimental drift, or rare events. A common approach is to train a classifier or anomaly detector and report performance metrics. However, good performance on one dataset does not by itself establish a general scientific finding. A more useful question is: which anomaly mechanisms are detectable by which feature groups, under which conditions, and why?

To address this, we design a controlled multi-domain experiment. Clean scientific signals are generated from six different domains: EEG-like signals, environmental sensor signals, spectroscopy-like signals, vibration signals, astronomical photometry-like signals, and fluid-sensor signals. Five controlled anomaly mechanisms are injected into these signals: spike, drift, dropout, oscillation disturbance, and distribution shift. Each anomaly is tested across multiple severity and noise levels. This gives a known ground truth for domain, anomaly mechanism, severity, and noise level.

The central claim tested in this paper is falsifiable:

\begin{quote}
This paper argues that SciPy-derived statistical and signal-processing features are effective for detecting controlled anomalies, but that detection performance alone does not imply diagnostic interpretability. In controlled multi-domain signals, binary detection was trivial, while anomaly-mechanism diagnosis remained difficult. We show that a major reason is structural: global energy features respond to many anomaly mechanisms, causing feature-group signatures to collapse toward amplitude-energy explanations. This demonstrates that scientific anomaly detection requires separating detection from diagnosis, and that mechanism diagnosis requires localized or mechanism-specific features rather than global summary statistics alone.
\end{quote}

The results support this claim in both controlled and external validation settings. SciPy-derived features make controlled anomalies easy to detect, but anomaly-mechanism classification is much harder, simple rule-based feature-group diagnosis performs at chance, and real-data validation shows that synthetic detection performance does not transfer directly to naturally occurring anomaly intervals.

The main contributions of this paper are:

\begin{itemize}
    \item a controlled multi-domain anomaly benchmark using six scientific signal domains and five anomaly mechanisms;
    \item a SciPy-based feature extraction workflow using statistical, signal-processing, spectral, residual, local-window, and complexity features;
    \item a comparison of SciPy-feature models, relative-signature models, and raw-signal baselines;
    \item an empirical demonstration that SciPy features strongly detect anomalies and preserve partial anomaly-mechanism information;
    \item a real-data validation experiment on the UCR Time Series Anomaly Archive;
    \item a negative result showing that simple dominant-feature-group explanations are insufficient for reliable anomaly diagnosis.
\end{itemize}

\section{Background and Motivation}

\subsection{Scientific Anomaly Detection}

An anomaly is an observation that differs from expected behavior. In scientific data, anomalies can represent measurement errors, sensor faults, experimental drift, unusual environmental conditions, or rare meaningful events. This makes anomaly detection scientifically delicate. Some anomalies should be removed as corrupted data, while others may be the most important observations in the dataset.

For this reason, anomaly detection in scientific workflows requires more than high predictive accuracy. It also requires diagnostic information. A scientist may need to distinguish between a spike caused by a sensor glitch, a drift caused by calibration change, a dropout caused by missing data, or a spectral disturbance caused by external interference. These cases may all appear as anomalies, but they imply different corrective actions.

\subsection{SciPy as an Anomaly Detection Toolkit}

SciPy provides many tools that are directly useful for anomaly detection. Statistical features from \texttt{scipy.stats} can summarize distributional shape, spread, skewness, kurtosis, entropy, and deviations from normality. Signal-processing features from \texttt{scipy.signal} can detect peaks, estimate local trends, compute residuals, and analyze spectral content. For example, the Savitzky Golay filter can be used to estimate a smooth local signal trend, while Welch's method estimates power spectral density by averaging periodograms over segments \citep{welch1967use}.

These features are attractive for scientific workflows because they are named, reproducible, computationally lightweight, and connected to interpretable signal properties. However, this does not automatically make them scientifically interpretable. A feature is not useful merely because it has a name. Interpretability must be operationalized: the feature or feature group should support a diagnostic conclusion about the anomaly mechanism.

\subsection{Detection Versus Diagnosis}

A central distinction in this paper is the difference between anomaly detection and anomaly diagnosis. Detection asks whether a signal is abnormal. Diagnosis asks what type of abnormality occurred. A model may perform very well at detection while still failing to identify the mechanism. This distinction is important because many anomaly-detection papers report high binary classification scores but do not show whether the model can distinguish between different anomaly causes.

This paper therefore evaluates both tasks. The binary task asks whether a signal is clean or anomalous. The multi-class task asks whether the anomaly mechanism is spike, drift, dropout, oscillation disturbance, or distribution shift. We also test a simple rule-based diagnosis system using relative feature-group signatures to evaluate whether intuitive feature groups can directly explain anomaly type.


\section{Related Work}

Time-series anomaly detection has a long history in statistical process monitoring. Classical approaches include threshold-based monitoring, Shewhart control charts, cumulative sum methods (CUSUM), and exponentially weighted moving average methods. These approaches are attractive because they are simple, interpretable, and computationally lightweight, but they often require assumptions about stationarity, distributional behavior, or known control limits.

Unsupervised machine-learning methods provide a more flexible alternative when labeled anomalies are rare. Isolation Forest is a widely used representative method that isolates anomalies through randomized partitioning rather than explicitly modeling normal density \citep{liu2008isolation}. One-class methods such as One-Class SVM similarly learn a boundary around normal data and are commonly used when only normal training samples are available. These methods are relevant to the UCR validation setting in this paper, where training regions are treated as normal and anomaly labels are used only for evaluation.

Deep learning methods have also been widely studied for time-series anomaly detection. LSTM encoder-decoder models learn to reconstruct normal temporal behavior and use reconstruction error as an anomaly score \citep{malhotra2016lstm}. More recent models such as THOC use temporal hierarchical one-class learning to capture multi-scale temporal structure \citep{shen2020thoc}, while transformer-based models such as TranAD use attention mechanisms for multivariate time-series anomaly detection and diagnosis \citep{tuli2022tranad}. These models can be powerful, but they are often more difficult to inspect and may require more data, tuning, and computational resources than classical feature-based methods.

Benchmark design is also a central issue in anomaly detection. Wu and Keogh argue that many time-series anomaly detection benchmarks are flawed and can create an illusion of progress, motivating more careful evaluation protocols and realistic anomaly archives \citep{wu2021ucr}. This concern is directly relevant to the present work: perfect detection on a controlled synthetic benchmark should not be interpreted as deployment-level success. For this reason, the controlled benchmark in this paper is used for mechanism testing, while the UCR archive is used as an external validation setting.

This paper differs from prior work by focusing on the distinction between anomaly detection and anomaly diagnosis. The goal is not to establish a new state-of-the-art detector. Instead, we test whether SciPy-derived statistical and signal-processing features provide a transparent representation that can detect anomalies, preserve partial mechanism-level information, and expose the limitations of simple feature-group interpretability.

\section{Methods}

\subsection{Overview}

The experiment consists of four steps. First, clean time-series signals are generated from multiple scientific domains. Second, controlled anomaly mechanisms are injected into those signals at different severity and noise levels. Third, SciPy statistical and signal-processing features are extracted from every signal. Fourth, models are trained and evaluated for binary anomaly detection, anomaly-type classification, relative signature analysis, and rule-based diagnosis.

Table~\ref{tab:learning-settings} summarizes the methodological difference between the controlled and UCR evaluations.

\begin{table}[htbp]
\centering
\caption{Learning settings used in the two evaluation stages.}
\label{tab:learning-settings}
\resizebox{\columnwidth}{!}{
\begin{tabular}{lll}
\toprule
\textbf{Experiment} & \textbf{Training Signal} & \textbf{Evaluation Goal} \\
\midrule
Controlled benchmark & Labeled anomalies & Detection and diagnosis \\
UCR validation & Normal windows only & Real-data detection \\
\bottomrule
\end{tabular}
}
\end{table}

\subsection{Scientific Signal Domains}

We generate clean signals from six scientific signal domains:

\begin{itemize}
    \item \textbf{EEG-like signals}: mixtures of low-frequency rhythms, burst-like oscillations, and colored noise;
    \item \textbf{Environmental sensor signals}: slow periodic cycles, low-frequency trends, and correlated noise;
    \item \textbf{Spectroscopy-like signals}: smooth baselines with Gaussian absorption or emission peaks;
    \item \textbf{Vibration signals}: oscillatory signals with harmonics and amplitude modulation;
    \item \textbf{Astronomical photometry-like signals}: light-curve-like signals with periodic dips and correlated noise;
    \item \textbf{Fluid-sensor signals}: quasi-periodic vortex-shedding-like signals with modulation and noise.
\end{itemize}

These are synthetic domain-inspired signals, not measurements from real instruments. They are used because controlled generation allows the true anomaly mechanism, severity, and noise level to be known. These domains were chosen because they differ in baseline structure, noise statistics, and signal dynamics. This allows the experiment to test whether anomaly behavior generalizes across heterogeneous scientific signals.

\subsection{Controlled Anomaly Mechanisms}

Five anomaly mechanisms are injected into the clean signals:

\begin{itemize}
    \item \textbf{Spike}: a short transient pulse added to the signal;
    \item \textbf{Drift}: a gradual baseline shift over part of the signal;
    \item \textbf{Dropout}: a flatline or attenuated region;
    \item \textbf{Oscillation disturbance}: a local high-frequency disturbance;
    \item \textbf{Distribution shift}: a local mean, variance, or scale change.
\end{itemize}

Each anomaly is evaluated across seven severity levels:
\[
0.02,\; 0.05,\; 0.08,\; 0.10,\; 0.15,\; 0.20,\; 0.35
\]
and four noise levels:
\[
1.0,\; 2.0,\; 3.0,\; 5.0.
\]
For each setting, 20 repeats are generated. Additional clean samples are generated to reduce class imbalance. The final controlled dataset contains 33,600 signals, each with 512 time points. The extracted feature matrix contains 56 SciPy-derived features.

\subsection{SciPy Feature Extraction}

For each signal \(x = [x_1, x_2, \ldots, x_T]\), we extract statistical, signal-processing, spectral, residual, local-window, and complexity features.

\subsubsection{Statistical Features}

The statistical feature group includes:

\begin{itemize}
    \item mean, standard deviation, variance;
    \item minimum, maximum, and range;
    \item median, interquartile range, and median absolute deviation;
    \item skewness and kurtosis;
    \item histogram entropy;
    \item normality-test statistics.
\end{itemize}

These features measure distributional shape, spread, central tendency, and tail behavior.

\subsubsection{Amplitude and Energy Features}

The amplitude-energy group includes:

\begin{itemize}
    \item absolute mean;
    \item root mean square amplitude;
    \item total energy;
    \item line length.
\end{itemize}

These features measure global signal magnitude and variation.

\subsubsection{Peak and Transient Features}

Peak and transient features include:

\begin{itemize}
    \item peak count;
    \item peak density;
    \item mean peak prominence;
    \item maximum peak prominence;
    \item mean peak width;
    \item maximum first derivative;
    \item maximum second derivative.
\end{itemize}

These features are intended to capture sharp transient events such as spikes.

\subsubsection{Trend and Residual Features}

Trend-residual features are computed using Savitzky Golay smoothing. The signal is smoothed, and residuals are computed as the difference between the original signal and the smoothed trend. Features include:

\begin{itemize}
    \item residual energy;
    \item residual absolute mean;
    \item maximum residual;
    \item maximum rolling mean change.
\end{itemize}

These features are intended to capture deviations from local trend and gradual drift.

\subsubsection{Dropout and Flatness Features}

Dropout-related features are computed from rolling variance:

\begin{itemize}
    \item minimum rolling variance;
    \item rolling variance ratio;
    \item flatness fraction.
\end{itemize}

These features are intended to capture flatline or attenuated regions.

\subsubsection{Spectral Features}

Spectral features are computed using Welch power spectral density estimates. They include:

\begin{itemize}
    \item bandpower values;
    \item bandpower ratios;
    \item spectral entropy;
    \item dominant frequency;
    \item high-frequency power ratio;
    \item local high-frequency ratio;
    \item local spectral entropy.
\end{itemize}

These features are intended to capture oscillation disturbances and changes in frequency structure.

\subsubsection{Complexity Features}

Complexity features include:

\begin{itemize}
    \item zero-crossing rate;
    \item Hjorth activity;
    \item Hjorth mobility;
    \item Hjorth complexity.
\end{itemize}

These features summarize temporal variability and signal complexity.

\subsection{Relative Feature-Group Signatures}

To test interpretability, features are grouped into seven feature groups:

\begin{itemize}
    \item statistical distribution;
    \item amplitude-energy;
    \item peak-transient;
    \item trend-residual;
    \item dropout-flatness;
    \item spectral-frequency;
    \item complexity.
\end{itemize}

For each domain, clean signals define a robust baseline. For each feature, we compute a robust absolute z-score relative to the clean baseline of the same domain. Extreme z-scores are clipped to reduce domination by large-magnitude features. Within each feature group, the median clipped z-score is used as the group response. The group responses are then normalized so that they sum to one. This produces a relative feature-group signature for each sample.

The goal is to test whether anomaly mechanisms map to distinct relative feature-group patterns.

\subsection{Models and Baselines}

We evaluate three representations:

\begin{itemize}
    \item full SciPy feature vectors;
    \item compressed relative feature-group signatures;
    \item raw time-domain signals.
\end{itemize}

For binary anomaly detection, we evaluate Logistic Regression, SVM-RBF, Random Forest, and Extra Trees models where applicable. For anomaly-type classification, we compare SciPy-feature models, relative-signature models, and raw-signal models.

The main evaluation metrics are balanced accuracy, precision, recall, F1-score, ROC-AUC, and PR-AUC for binary anomaly detection. For anomaly-type classification, we report accuracy, balanced accuracy, macro-F1, and weighted-F1.

\subsection{Rule-Based Diagnosis}

To test whether feature groups directly support mechanistic diagnosis, we define a simple rule-based diagnosis system:

\begin{itemize}
    \item spike: peak-transient, statistical-distribution, and trend-residual groups;
    \item drift: trend-residual and statistical-distribution groups;
    \item dropout: dropout-flatness and trend-residual groups;
    \item oscillation disturbance: spectral-frequency and complexity groups;
    \item distribution shift: statistical-distribution and amplitude-energy groups.
\end{itemize}

For each anomalous signal, the rule system computes the mean response of the expected groups for each anomaly type and predicts the type with the largest score. If simple feature-group interpretability is sufficient, this rule system should classify anomaly type above chance.

\section{Results}

\subsection{Dataset Summary}

The final controlled experiment produced 33,600 signals with 512 time points each. The feature extraction pipeline produced 56 SciPy-derived features per signal. The experiment included six domains, five anomaly mechanisms, seven severity levels, four noise levels, and 20 anomaly repeats per setting.

\subsection{Binary Detection Was Trivial Under Controlled Injection}

Binary anomaly detection was nearly solved by the SciPy feature representation. Random Forest trained on SciPy features achieved balanced accuracy, F1-score, ROC-AUC, and PR-AUC of 1.0000. Relative Signatures + Logistic Regression also achieved perfect binary detection. Although this demonstrates that the injected anomalies are measurable in SciPy feature space, it also shows that binary detection is not the scientifically interesting part of the experiment. A perfect detection score indicates that the clean/anomalous boundary is too easily separable under the current synthetic design.

For this reason, we do not interpret the binary detection result as evidence of a strong anomaly detector. Instead, we treat it as evidence that the controlled perturbations produced measurable feature deviations. The central question is whether those deviations preserve information about anomaly mechanism. Table~\ref{tab:binary-detection} reports the binary detection results.

\begin{table*}[htbp]
\centering
\caption{Binary anomaly detection results.}
\label{tab:binary-detection}
\resizebox{\textwidth}{!}{
\begin{tabular}{lccccc}
\toprule
\textbf{Model} & \textbf{Accuracy} & \textbf{Balanced Acc.} & \textbf{F1} & \textbf{ROC-AUC} & \textbf{PR-AUC} \\
\midrule
SciPy Features + Random Forest & \textbf{1.0000} & \textbf{1.0000} & \textbf{1.0000} & \textbf{1.0000} & \textbf{1.0000} \\
Relative Signatures + Logistic Regression & \textbf{1.0000} & \textbf{1.0000} & \textbf{1.0000} & \textbf{1.0000} & \textbf{1.0000} \\
SciPy Features + Extra Trees & 0.9995 & 0.9995 & 0.9995 & 1.0000 & 1.0000 \\
SciPy Features + SVM-RBF & 0.9190 & 0.8490 & 0.9568 & 0.9378 & 0.9761 \\
SciPy Features + Logistic Regression & 0.7975 & 0.7825 & 0.8846 & 0.8729 & 0.9279 \\
Raw Signal + Extra Trees & 0.7060 & 0.7060 & 0.6632 & 0.7993 & 0.8198 \\
Raw Signal + Random Forest & 0.7073 & 0.7073 & 0.6650 & 0.7926 & 0.8166 \\
Raw Signal + SVM-RBF & 0.9340 & 0.5411 & 0.9658 & 0.6177 & 0.8082 \\
Raw Signal + Logistic Regression & 0.7109 & 0.4708 & 0.8298 & 0.4747 & 0.5295 \\
\bottomrule
\end{tabular}
}
\end{table*}

The raw-signal tree baselines performed substantially worse than the SciPy-feature models, with balanced accuracy around 0.71 and F1 around 0.66. This shows that raw time-domain samples contain some anomaly information, but the SciPy feature representation makes the clean/anomalous boundary much more separable.

\subsection{Anomaly Diagnosis Was Substantially Harder}

The anomaly-type classification task was much harder than binary detection. The best model, Random Forest trained on the full SciPy feature set, achieved macro-F1 = 0.4578 across five anomaly mechanisms. This is well above the raw-signal Extra Trees baseline, which achieved macro-F1 = 0.2051, approximately chance for five classes. Because there are five anomaly mechanisms, chance-level macro-F1 is approximately 0.20 under balanced classes.

This result shows that SciPy features preserve some mechanism-level information, but not enough for robust diagnosis across all anomaly types and domains. In other words, SciPy features can detect that a signal has changed, but determining how it changed is a harder problem. Table~\ref{tab:type-classification} reports the anomaly diagnosis results.

\begin{table*}[htbp]
\centering
\caption{Anomaly-type classification results across five anomaly mechanisms.}
\label{tab:type-classification}
\begin{tabular}{lccc}
\toprule
\textbf{Model} & \textbf{Accuracy} & \textbf{Balanced Acc.} & \textbf{Macro-F1} \\
\midrule
SciPy Features + Random Forest & \textbf{0.4705} & \textbf{0.4705} & \textbf{0.4578} \\
SciPy Features + Extra Trees & 0.4090 & 0.4090 & 0.4122 \\
Relative Signatures + Extra Trees & 0.2148 & 0.2148 & 0.2149 \\
Relative Signatures + Logistic Regression & 0.2183 & 0.2183 & 0.2064 \\
Raw Signal + Extra Trees & 0.2045 & 0.2045 & 0.2051 \\
\bottomrule
\end{tabular}
\end{table*}

\subsection{Confusion Structure of Anomaly Mechanisms}

The macro-F1 score shows that anomaly-mechanism classification is substantially harder than binary detection, but it does not reveal which mechanisms are confused. We therefore inspect the confusion matrix of the best anomaly-type classifier, SciPy Features + Random Forest. Table~\ref{tab:type-confusion-normalized} shows the row-normalized confusion matrix, where each row represents the true anomaly type and each column represents the predicted anomaly type.

\begin{table*}[htbp]
\centering
\caption{Row-normalized confusion matrix for anomaly-type classification using SciPy Features + Random Forest. Rows indicate true anomaly type and columns indicate predicted anomaly type.}
\label{tab:type-confusion-normalized}
\resizebox{\textwidth}{!}{
\begin{tabular}{lccccc}
\toprule
\textbf{True / Predicted} & \textbf{Distribution Shift} & \textbf{Drift} & \textbf{Dropout} & \textbf{Oscillation Disturbance} & \textbf{Spike} \\
\midrule
\textbf{Distribution Shift} & 0.34 & 0.17 & 0.07 & 0.28 & 0.13 \\
\textbf{Drift} & 0.07 & 0.72 & 0.03 & 0.12 & 0.06 \\
\textbf{Dropout} & 0.06 & 0.09 & 0.60 & 0.10 & 0.15 \\
\textbf{Oscillation Disturbance} & 0.21 & 0.04 & 0.11 & 0.45 & 0.18 \\
\textbf{Spike} & 0.24 & 0.10 & 0.18 & 0.26 & 0.21 \\
\bottomrule
\end{tabular}
}
\end{table*}

The confusion matrix shows that the classifier learned some anomaly mechanisms better than others. Drift was the most separable mechanism, with 72\% of drift samples correctly classified. Dropout was also partially separable, with 60\% correctly classified. Oscillation disturbance was detected at a lower but still meaningful rate, with 45\% correctly classified. In contrast, distribution shift and spike were much less separable. Distribution shift was correctly classified only 34\% of the time and was often confused with oscillation disturbance. Spike was the weakest class, with only 21\% correctly classified, and was frequently confused with distribution shift and oscillation disturbance.

This confusion structure supports the paper's central argument: SciPy features contain partial mechanism-level information, but anomaly diagnosis is not solved by global feature summaries. Drift and dropout produce more consistent feature patterns, likely because they alter baseline or flatness structure. Spike and oscillation disturbance are more ambiguous in the current representation, suggesting that short local transients and local spectral changes are not sufficiently isolated by whole-signal features. This reinforces the need for local anomaly localization before mechanism diagnosis.

This pattern has a mechanistic explanation. Drift produces a sustained baseline displacement, which changes rolling mean structure and creates a cumulative low-frequency effect across a large portion of the signal. Because the perturbation persists over many time steps, global and rolling-window features can capture it reliably. Dropout similarly creates an extended low-variance or attenuated region, which is partly captured by rolling-variance and flatness features. Spike is different: it is a short local transient. When features are computed over the full \(T=512\) point segment, the contribution of a narrow spike is diluted by the segment length. A local spike may strongly affect a few derivative or peak features, but many global statistics average it with hundreds of normal samples. This explains why spike was frequently confused with distribution shift and oscillation disturbance. It also suggests that reliable spike diagnosis requires anomaly localization or local-window feature extraction before whole-signal classification.

\subsection{Relative Feature-Group Signature Map}

Table~\ref{tab:signature-map} shows the average relative feature-group signature for each anomaly type. The strongest group for every anomaly type was amplitude-energy.

\begin{table*}[htbp]
\centering
\caption{Relative feature-group signature map by anomaly type.}
\label{tab:signature-map}
\begin{tabular}{lccccccc}
\toprule
\textbf{Anomaly Type} & \textbf{Stat.} & \textbf{Amp./Energy} & \textbf{Peak} & \textbf{Trend} & \textbf{Dropout} & \textbf{Spectral} & \textbf{Complexity} \\
\midrule
Distribution shift & 0.1046 & \textbf{0.5383} & 0.0745 & 0.0855 & 0.0572 & 0.0598 & 0.0801 \\
Drift & 0.1069 & \textbf{0.5379} & 0.0733 & 0.0835 & 0.0583 & 0.0599 & 0.0803 \\
Dropout & 0.1088 & \textbf{0.5279} & 0.0756 & 0.0833 & 0.0663 & 0.0612 & 0.0770 \\
Oscillation disturbance & 0.1044 & \textbf{0.5403} & 0.0735 & 0.0844 & 0.0583 & 0.0599 & 0.0792 \\
Spike & 0.1046 & \textbf{0.5407} & 0.0738 & 0.0846 & 0.0566 & 0.0602 & 0.0796 \\
\bottomrule
\end{tabular}
\end{table*}

The signature map does not support the simple hypothesis that each anomaly type maps cleanly to a different dominant feature group. Even after robust normalization, clipping, median aggregation, and relative normalization, amplitude-energy remained the dominant group for every anomaly type. This indicates that the controlled anomaly injections often changed global signal energy even when the intended anomaly mechanism was local, spectral, or structural.

\subsection{Rule-Based Diagnosis}

The rule-based diagnosis system achieved overall accuracy = 0.1993, approximately chance for five classes. Table~\ref{tab:rule-diagnosis} shows the accuracy by anomaly type.

\begin{table}[htbp]
\centering
\caption{Rule-based diagnosis accuracy by anomaly type.}
\label{tab:rule-diagnosis}
\begin{tabular}{lc}
\toprule
\textbf{Anomaly Type} & \textbf{Rule Accuracy} \\
\midrule
Distribution shift & 0.9917 \\
Drift & 0.0006 \\
Dropout & 0.0063 \\
Oscillation disturbance & 0.0006 \\
Spike & 0.0000 \\
\bottomrule
\end{tabular}
\end{table}

The rule-based diagnosis result reveals the central failure mode of naive interpretability. Distribution shift achieved 0.9917 rule accuracy, while spike, drift, dropout, and oscillation disturbance were nearly always misdiagnosed. This pattern is not random. The rule system assigns distribution shift using statistical-distribution and amplitude-energy responses. Since amplitude-energy dominated the relative signature map for every anomaly type, the rule system was biased toward the distribution-shift rule. As a result, distribution shift was diagnosed correctly, not because the rule system learned a rich mechanism-specific signature, but because the confounding amplitude-energy response made other anomaly types resemble distribution-shift behavior.

This explains why the overall rule-based diagnosis accuracy was approximately chance. The rule system did not fail because the features were meaningless. It failed because a globally responsive feature group collapsed the signature space. This is a precise example of why named features alone are not sufficient for scientific interpretability.

\subsection{Leave-One-Domain-Out Generalization}

To test domain generalization, we trained an anomaly-type classifier on five domains and tested on the held-out sixth domain. Table~\ref{tab:logo} shows the results.

\begin{table}[htbp]
\centering
\caption{Leave-one-domain-out anomaly-type classification.}
\label{tab:logo}
\begin{tabular}{lcc}
\toprule
\textbf{Held-Out Domain} & \textbf{Balanced Acc.} & \textbf{Macro-F1} \\
\midrule
EEG-like & 0.4114 & 0.4194 \\
Fluid sensor & 0.3939 & 0.3941 \\
Astronomical photometry & 0.3800 & 0.3728 \\
Environmental sensor & 0.3418 & 0.3503 \\
Vibration & 0.3307 & 0.3291 \\
Spectroscopy & 0.2611 & 0.2087 \\
\bottomrule
\end{tabular}
\end{table}

The leave-one-domain-out results show partial but limited domain generalization. Some domains, such as EEG-like and fluid-sensor signals, generalized better than others. Spectroscopy and vibration were more difficult. This suggests that anomaly mechanisms are not fully domain-independent. The same injected anomaly may interact differently with different baseline signal structures.

\subsection{Statistical Tests on Feature-Group Responses}

Kruskal Wallis tests were used to test whether feature-group responses differed by anomaly type. Table~\ref{tab:kruskal} shows the results.

\begin{table}[htbp]
\centering
\caption{Kruskal Wallis tests for feature-group response differences by anomaly type.}
\label{tab:kruskal}
\begin{tabular}{lcc}
\toprule
\textbf{Feature Group} & \textbf{H Statistic} & \textbf{p-value} \\
\midrule
Statistical distribution & 56.2222 & \(1.80 \times 10^{-11}\) \\
Dropout flatness & 36.5159 & \(2.27 \times 10^{-7}\) \\
Amplitude-energy & 26.7147 & \(2.27 \times 10^{-5}\) \\
Complexity & 9.9531 & \(4.12 \times 10^{-2}\) \\
Spectral frequency & 6.3235 & 0.1763 \\
Trend residual & 6.2548 & 0.1809 \\
Peak transient & 5.4034 & 0.2484 \\
\bottomrule
\end{tabular}
\end{table}

The statistically significant differences appeared in statistical-distribution, dropout-flatness, amplitude-energy, and complexity groups. However, not all feature groups showed significant separation. Spectral-frequency, trend-residual, and peak-transient groups were not significant in this run. This again suggests that anomaly-type differences exist, but they are not captured by a simple one-group-per-anomaly interpretation.

An important irony is that the peak-transient and trend-residual groups, which are the most intuitively matched to spike and drift anomalies, did not reach Kruskal Wallis significance in this run. This reinforces the central argument that intuitive feature names do not automatically translate into reliable diagnostic separation.

\section{Real-Data Validation on the UCR Time Series Anomaly Archive}

\subsection{Dataset and Setup}

To evaluate whether the SciPy anomaly detection workflow transfers beyond controlled synthetic injection, we performed an external validation experiment on the UCR Time Series Anomaly Archive \citep{wu2021ucr}. Unlike the controlled benchmark, where anomaly type, severity, and noise level are generated by design, the UCR archive contains real time-series anomaly problems with labeled anomaly intervals encoded in the file metadata.

We evaluated the first 60 archive files that could be parsed by the experiment script. Four files were skipped because the sliding-window labeling procedure produced no anomalous test windows, leaving 56 evaluated files. Each time series was converted into overlapping windows of length 128 with step size 16. Windows ending before the training boundary were used as normal training windows. Test windows were labeled anomalous if they overlapped the labeled anomaly interval by at least 10\% of the window length. Point-adjusted F1 treats an anomaly interval as detected if at least one overlapping window is flagged, which is common in time-series anomaly detection but can be more optimistic than raw point-level F1.

For each window, we extracted the same SciPy-inspired statistical and signal-processing features used in the controlled experiment, including distributional statistics, peak features, residual features, rolling-window features, and spectral features. We evaluated four unsupervised methods: SciPy Features + Robust Distance, SciPy Features + Isolation Forest, SciPy Features + One-Class SVM, and Raw Windows + Isolation Forest. We do not claim state-of-the-art performance on the UCR archive. The purpose of the UCR experiment is external validation of the SciPy feature workflow under naturally labeled anomaly intervals, not benchmark optimization.

\paragraph{Methodological asymmetry.}
The controlled benchmark and UCR validation evaluate different learning settings. In the controlled benchmark, anomaly labels and anomaly-mechanism labels are known by construction, so supervised classifiers can be trained for both binary detection and anomaly-type classification. In contrast, the UCR archive provides normal training regions and labeled anomaly intervals for evaluation, but it does not provide labeled examples of each anomaly mechanism for supervised mechanism classification. Therefore, the UCR validation is evaluated using unsupervised and one-class methods trained on normal windows. This asymmetry is intentional: the controlled benchmark tests mechanism-level hypotheses under known labels, while the UCR experiment tests whether the SciPy feature workflow remains useful when only normal training data are available. The two experiments should therefore not be interpreted as identical evaluations, but as complementary tests of mechanism understanding and real-data transfer.

\subsection{Results}

Table~\ref{tab:ucr-results} summarizes the real-data validation results. The strongest method by mean point-adjusted F1 was SciPy Features + Isolation Forest, which achieved mean point-adjusted F1 = 0.2884, raw F1 = 0.1248, ROC-AUC = 0.7758, PR-AUC = 0.1801, and event recall = 0.4107 across 56 evaluated files. SciPy Features + One-Class SVM achieved the highest event recall at 0.6607, but had a lower mean point-adjusted F1 of 0.2646.

\begin{table*}[htbp]
\centering
\caption{Real-data validation on the UCR Time Series Anomaly Archive.}
\label{tab:ucr-results}
\begin{tabular}{lcccccc}
\toprule
\textbf{Method} & \textbf{Raw Bal. Acc.} & \textbf{Raw F1} & \textbf{ROC-AUC} & \textbf{PR-AUC} & \textbf{Point-Adj. F1} & \textbf{Event Recall} \\
\midrule
SciPy Features + Isolation Forest & 0.5588 & 0.1248 & \textbf{0.7758} & 0.1801 & \textbf{0.2884} & 0.4107 \\
SciPy Features + One-Class SVM & \textbf{0.6374} & \textbf{0.1375} & 0.6561 & \textbf{0.2143} & 0.2646 & \textbf{0.6607} \\
Raw Windows + Isolation Forest & 0.5484 & 0.1184 & 0.5450 & 0.1582 & 0.2617 & 0.3571 \\
SciPy Features + Robust Distance & 0.5340 & 0.0877 & 0.7225 & 0.1331 & 0.2328 & 0.2857 \\
\bottomrule
\end{tabular}
\end{table*}

The UCR results show a sharp contrast with the controlled synthetic benchmark. In the synthetic experiment, binary anomaly detection was nearly trivial, with SciPy-feature models reaching perfect performance. In the real-data validation experiment, performance was much lower. This confirms that perfect performance under controlled injection should not be interpreted as deployment-level anomaly detection ability. Real anomaly intervals are more heterogeneous, less aligned with fixed window boundaries, and often subtler than injected perturbations.

The real-data results also show that different metrics emphasize different operational goals. SciPy Features + Isolation Forest achieved the best point-adjusted F1 and ROC-AUC, suggesting better ranking and localization performance. SciPy Features + One-Class SVM achieved the highest event recall, suggesting that it was more likely to flag at least one window inside the anomaly interval. Depending on the scientific workflow, event recall may be more important than point-level precision if the goal is to alert a researcher that a suspicious episode occurred.

\subsection{Best and Worst Cases}

The SciPy Features + Isolation Forest method performed strongly on some UCR files. For example, it achieved point-adjusted F1 values above 0.90 on selected BIDMC, CIMIS air temperature, PowerDemand, and Lab2Cmac files. However, it failed completely on several PowerDemand, TkeepMARS, apnea ECG, and gait files, where point-adjusted F1 was 0.0 and the anomaly event was not detected.

This variation is important. It shows that SciPy feature-based anomaly detection is not universally reliable across real time-series domains. Some anomaly intervals produce feature deviations that are detectable by the current window-level feature representation, while others do not. This supports the paper's broader conclusion: SciPy features are useful and interpretable, but real scientific anomaly detection requires domain calibration, careful thresholding, and possibly local or task-specific feature design.

\section{Discussion}

\subsection{Detection Is Easier Than Diagnosis}

The most important finding is the gap between anomaly detection and anomaly diagnosis. Detecting whether a signal is abnormal was easy for SciPy-derived features. Both full SciPy features and compact relative signatures achieved perfect binary detection. However, identifying the specific anomaly mechanism was much harder. The best SciPy-feature model achieved macro-F1 = 0.4578, while raw-signal features were near chance.

This distinction matters because many anomaly-detection systems report strong binary detection scores but provide little information about the nature of the anomaly. In scientific workflows, a detection flag is often only the first step. The more useful question is whether the system can help the scientist understand what happened.

\subsection{Synthetic Performance Does Not Transfer Directly to Real Data}

The real-data validation experiment clarifies the role of the controlled synthetic benchmark. In the synthetic experiment, binary detection was perfect, which indicates that the injected anomalies produced measurable feature deviations. However, on the UCR Time Series Anomaly Archive \citep{wu2021ucr}, the best method achieved only 0.2884 mean point-adjusted F1 and 0.4107 event recall. This gap shows that controlled injection experiments should be used to study mechanisms, not to claim deployment-level anomaly-detection performance.

The real-data results strengthen the central argument of this paper. Detection and diagnosis are different problems, and both become more difficult when anomaly structure is not controlled by the experimenter. The UCR archive introduces heterogeneous anomaly shapes, different baseline dynamics, and difficult thresholding conditions. Under these conditions, SciPy features remain useful, especially in ROC-AUC ranking, but they do not solve anomaly detection. This supports a cautious interpretation: SciPy provides a strong and transparent baseline for scientific anomaly workflows, but it should be paired with domain-specific validation.

\subsection{Naive Interpretability Failed}

The rule-based diagnosis system was intentionally simple. It tested whether intuitive feature groups could directly diagnose anomaly mechanism. For example, spikes were expected to activate peak-transient features, drift was expected to activate trend-residual features, dropout was expected to activate flatness features, and oscillation disturbance was expected to activate spectral features.

The results did not support this simple interpretation. The rule-based system achieved chance-level accuracy. The average relative signature map was dominated by amplitude-energy for every anomaly type. This shows that named features alone are not enough for scientific interpretability. A feature can be interpretable in isolation but still fail to provide a reliable diagnostic explanation when used in a complex multi-domain setting.

\subsection{Why Global Energy Features Confound Diagnosis}

The dominance of amplitude-energy features can be understood directly from the structure of additive anomalies. Let \(x \in \mathbb{R}^T\) be a clean signal and let an anomaly be represented as a perturbation \(\delta \in \mathbb{R}^T\). The observed anomalous signal is

\[
\tilde{x} = x + \delta .
\]

The squared energy of the anomalous signal is

\[
\|\tilde{x}\|_2^2 = \|x + \delta\|_2^2 .
\]

Expanding this expression gives

\[
\|\tilde{x}\|_2^2
=
\|x\|_2^2
+
2\langle x,\delta \rangle
+
\|\delta\|_2^2 .
\]

Therefore, the signed change in energy is

\[
\Delta E
=
\|\tilde{x}\|_2^2 - \|x\|_2^2
=
2\langle x,\delta \rangle + \|\delta\|_2^2 .
\]

For additive spikes, oscillatory disturbances, and many distribution shifts, this quantity often increases. However, not all anomalies increase energy. A dropout or attenuation event can be represented locally by \(\delta \approx -x\), which may reduce signal energy in the affected region. Therefore, the relevant diagnostic issue is not only the sign of \(\Delta E\), but the magnitude of the deviation from the clean-domain energy baseline:

\[
|\Delta E|
=
\left|
\|\tilde{x}\|_2^2 - \|x\|_2^2
\right|.
\]

Amplitude-energy features such as RMS amplitude, absolute mean, variance, and total energy can respond to both increases and decreases in signal magnitude. This makes them mechanism-agnostic responders: they can detect that the signal's magnitude structure changed without identifying whether the cause was a spike, dropout, drift, oscillation disturbance, or distribution shift.

This explains why amplitude-energy dominated the relative feature-group signature map for all anomaly types. Spike, drift, dropout, oscillation disturbance, and distribution shift all modify signal magnitude structure in some way. As a result, global energy features act as a confounder for mechanism diagnosis. They are useful for detection but structurally insufficient for distinguishing anomaly type.

\subsection{SciPy Features Still Preserve Mechanism Information}

Although simple feature-group rules failed, full SciPy features did preserve partial mechanism information. The best SciPy-feature classifier achieved macro-F1 = 0.4578, compared with 0.2051 for the raw-signal baseline. This means SciPy features contain useful information about anomaly type, but the information is multivariate and overlapping.

This result supports a more realistic interpretation of scientific anomaly diagnosis. Mechanisms may not correspond to one feature group. Instead, they may appear as distributed patterns across statistical, amplitude, residual, spectral, flatness, peak, and complexity features.

\subsection{Domain Generalization Is Partial}

The leave-one-domain-out experiment showed that anomaly-type classification generalizes only partially across domains. EEG-like and fluid-sensor domains were easier to generalize to than spectroscopy and vibration. This suggests that domain dynamics matter. Spectroscopy was especially difficult because the clean spectroscopy-like baseline already contains peak-like structures, making spike-like anomalies harder to separate from normal signal morphology. Vibration was also difficult because oscillatory behavior is already part of the normal baseline, so oscillation disturbance can look like ordinary signal variation.

Therefore, the paper's claim should not be that SciPy features produce universal anomaly signatures. A more defensible claim is that SciPy features provide a useful representation for anomaly detection and partial mechanism classification, while domain-specific calibration remains necessary.


\section{Reproducibility}

All experiments were implemented in Python using NumPy, SciPy, scikit-learn, pandas, and Matplotlib. A fixed random seed of 42 was used for controlled signal generation, train-test splitting, and model initialization where supported. Each controlled signal contained 512 time points. The controlled experiment used six generated scientific signal domains, five injected anomaly mechanisms, seven severity levels, four noise levels, and 20 anomaly repeats per setting. Additional clean samples were generated to reduce class imbalance. The UCR validation used sliding windows of length 128 with step size 16, and test windows were labeled anomalous if at least 10\% of the window overlapped the labeled anomaly interval.

\section{Limitations}

This study has several limitations. First, the main controlled benchmark uses synthetic anomaly injection rather than naturally labeled instrument failures. This is useful for mechanistic testing because the true anomaly type, severity, and noise level are known, but the UCR validation shows that synthetic performance does not directly imply real-data robustness.

Second, the binary anomaly-detection task is still too easy in the current benchmark. Even at low severity and high noise, SciPy-feature models achieved perfect detection. Future work should use weaker anomalies, more realistic nuisance variation, and harder clean/anomaly overlap.

Third, the relative feature-group signatures were computed from global signal summaries. This may be insufficient for local anomaly mechanisms. Future versions should first localize anomalous regions and then compute local diagnostic features.

Fourth, the rule-based diagnosis system was simple. Its failure does not prove that interpretable diagnosis is impossible. It shows that naive dominant-feature-group explanations are insufficient.

Fifth, the generated domains are synthetic approximations of scientific signals. They are useful for controlled testing, but real scientific datasets may contain more complex noise, artifacts, missingness, and domain-specific structure. The UCR archive provides one external validation setting, but it does not replace domain-specific validation on the scientific instruments and workflows where the detector would actually be used.

Another threat to validity is that the synthetic anomaly mechanisms are simplified mathematical perturbations. Although they allow controlled mechanism testing, real faults may combine multiple mechanisms at once, such as drift plus dropout or oscillation plus amplitude scaling.

Another limitation is that the controlled and UCR experiments use different learning settings: supervised classification is possible in the controlled benchmark because anomaly labels are generated by design, while the UCR validation uses unsupervised one-class methods because only normal training windows are available. Therefore, the numerical results from the two experiments should not be compared directly.

\section{Future Work}

Future work should extend this study in several directions. First, external validation should be expanded beyond the UCR archive to real fault-labeled datasets from environmental monitoring, spectroscopy, astronomical photometry, industrial vibration, and laboratory sensors. Second, anomaly localization should be added so that local features can be computed within the anomalous region rather than over the full signal.

A concrete path forward is to compute diagnostic features after anomaly localization rather than over the entire \(T=512\) signal. For example, spike diagnosis should be based on a local window centered around the detected transient, where peak prominence, derivative magnitude, and residual deviation are not diluted by hundreds of normal time points. Similarly, drift diagnosis should use features computed over the affected interval, such as rolling-mean slope, cumulative deviation, and local low-frequency residual structure. This would turn the current global feature pipeline into a two-stage system: first detect and localize the anomalous region, then compute mechanism-specific SciPy features within that region for diagnosis.

Third, more realistic anomaly injections should be developed, including overlapping anomalies and nuisance shifts that are not true anomalies. Fourth, interpretability should be evaluated not only by feature names but by whether the explanation helps a user identify the correct anomaly mechanism.

Future work should also compare SciPy-feature models with deep-learning models under limited-data conditions. A key practical question is whether SciPy features can approach deep-learning performance while being cheaper, easier to inspect, and more data-efficient.

\section{Conclusion}

The controlled and real-data experiments together support a cautious but useful conclusion. SciPy statistical and signal-processing features make controlled anomalies easy to detect, but high detection performance under synthetic injection does not imply real-world robustness or diagnostic interpretability. In the controlled benchmark, anomaly-mechanism classification remained difficult and simple rule-based feature-group diagnosis performed at chance. In the UCR real-data validation, the best SciPy method achieved only 0.2884 mean point-adjusted F1, showing that although SciPy features perfectly separate synthetic clean/anomalous signals, real anomaly intervals remain difficult. These findings suggest that SciPy features are best understood as a transparent baseline and diagnostic starting point, not as a complete anomaly-detection solution. Scientific anomaly detection should distinguish between detecting that a signal is abnormal, diagnosing why it is abnormal, and validating whether the method transfers to naturally occurring anomalies.
